\section{A fully spectral kinematic match}
\label{sec:spectral-km}

\hspace*{1em}We derive a model of droplet--bath contact in which the bath and droplet are represented as coupled spectral oscillator systems, following \cite{AlventosaEtAl2023}; contact is enforced over the entire pressed region and the pressure is solved for directly, following \cite{AgueroEtAl2026}; and the map between the bath's cylindrical coordinate and the droplet's polar angle is used only in its explicit, forward direction. The section is organized as follows: \S\ref{subsec:formulation} fixes notation and the geometric relation between the bath's and droplet's coordinates; \S\ref{subsec:bath} and \S\ref{subsec:drop} state the (unmodified) bath and droplet models; \S\ref{subsec:affine} derives, once, the affine dependence of each interface's state on the pressure induced by implicit time-stepping; \S\ref{subsec:contact} derives the contact closure; \S\ref{subsec:newton} solves it as a one-dimensional outer iteration over the contact angle wrapping an inner solve for the pressure; \S\ref{subsec:summary} collects the complete, self-contained system solved at each time step; and \S\ref{subsec:corrections} records arguments from earlier revisions of this document that were wrong, together with the structure removed along with them.

\subsection{Governing equations and interface parametrization}
\label{subsec:formulation}

We use the nondimensionalization of \cite{AlventosaEtAl2023}: length scale $R$ (the undeformed droplet radius), time scale $t_\sigma=\sqrt{\rho R^3/\sigma}$, force scale $2\pi\sigma R$, giving the Weber, Bond and Ohnesorge numbers $\We,\Bo,\Oh$ as defined there. Cylindrical coordinates $(r,z)$ describe the bath, with $z=0$ the undisturbed free surface and gravity acting in $-\hat z$; spherical coordinates $(\xi,\theta)$, sharing the same axis, are centred on the droplet's centre of mass, with $\theta=0$ along $-\hat z$ so that $\theta=0$ is the point of the droplet nearest the bath. The bath interface height and the droplet radius are decomposed as
\begin{equation}
    \eta(r,\tau) = \sum_{m=0}^{M} a_m(\tau)\,J_0(k_mr)\text{,}\qquad
    \xi(\theta,\tau) = 1+\sum_{l=2}^{L}\beta_l(\tau)\,P_l(\cos\theta)\text{,}
    \label{eq:eta-xi}
\end{equation}
and the droplet's centre of mass is at height $z_{cm}(\tau)$ above the undisturbed bath surface. A point on the droplet surface at polar angle $\theta$ lies, in the bath's cylindrical coordinates, at horizontal distance and height
\begin{equation}
    r(\theta,\tau) = \xi(\theta,\tau)\sin\theta\text{,}\qquad
    z_d(\theta,\tau) = \xi(\theta,\tau)\cos\theta
    \label{eq:forward-map}
\end{equation}
below the centre of mass (so a point's absolute height, in the bath's $z$-coordinate, is $z_{cm}(\tau)-z_d(\theta,\tau)$: $z_d(0,\tau)=\xi(0,\tau)>0$ places $\theta=0$ below $z_{cm}$, consistent with it being the point nearest the bath), so that the droplet's underside, at the same instant, meets the bath surface if and only if $\eta(r(\theta,\tau),\tau)=z_{cm}(\tau)-z_d(\theta,\tau)$. Equation~\eqref{eq:forward-map} is a map from $\theta$ to $(r,z_d)$, given explicitly, for any $\theta$, once $\{\beta_l(\tau)\}$ is known; it is never inverted for $\theta$ given $r$ in what follows. We write $x:=\cos\theta$ throughout when convenient, so that $P_l(\cos\theta)=P_l(x)$.

\subsection{Bath interface model}
\label{subsec:bath}

The bath is governed by the linearized, weakly viscous free-surface equations of \cite{AlventosaEtAl2023}, themselves following \cite{DiasEtAl2008}: with velocity potential $\phi(r,z,\tau)$,
\begin{equation}
    \nabla_H^2\phi+\partial_z^2\phi=0\text{,}\quad z\le0\text{;}\qquad
    \partial_\tau\eta = \partial_z\phi+2\Oh\nabla^2\eta\text{,}\quad z=0\text{;}
    \label{eq:bath-pde-1}
\end{equation}
\begin{equation}
    \partial_\tau\phi = -\Bo\,\eta - 2\Oh\,\partial_z^2\phi + \kappa[\eta] - p\text{,}\quad z=0\text{;}\qquad
    \partial_z\phi=0\text{ at }z=-h_0\text{,}
    \label{eq:bath-pde-2}
\end{equation}
with $\nabla_H^2=\partial_r^2+r^{-1}\partial_r$ the horizontal Laplacian (written $\nabla^2$ in the surface conditions above, where only the horizontal part acts) and $\kappa[\eta]=\nabla_H^2\eta$ the linearized curvature and $h_0$ the (nondimensional) bath depth (at $z=-h_0$, consistent with $z\le0$; \citealt{AlventosaEtAl2023} print $z=h_0$), subject to the no-flux wall condition $\partial_r\eta=\partial_r\phi=0$ at $r=b$. Expanding $\eta$ as in \eqref{eq:eta-xi} on the eigenfunctions of $(\nabla^2+k_m^2)J_0(k_mr)=0$ satisfying $J_0'(k_mb)=0$, and eliminating $\phi$ between \eqref{eq:bath-pde-1} and \eqref{eq:bath-pde-2} mode by mode, gives one damped linear oscillator per bath mode,
\begin{equation}
    \ddot a_m + 4\Oh k_m^2 \dot a_m + \left(k_m^2+\Bo\right)k_m\tanh(k_m h_0)\, a_m = -2 k_m \tanh(k_m h_0)\, c_m(\tau)\text{,}\qquad m=0,\dots,M\text{,}
    \label{eq:bath-mode}
\end{equation}
where
\begin{equation}
    c_m(\tau) = \frac{2}{(bJ_0(k_mb))^2}\int_0^b p(r,\tau)\,r\,J_0(k_mr)\,dr
    \label{eq:c_m-def}
\end{equation}
is the $m$-th Fourier--Bessel coefficient of the actual contact pressure $p$, extended by zero outside the pressed region. The normalization deserves a remark, because it differs from the one printed by \cite{AlventosaEtAl2023} and the difference is not cosmetic. The no-flux wall condition gives $J_0'(k_mb)=0$, equivalently $J_1(k_mb)=0$, so $k_mb$ is the $m$-th zero of $J_1$ and the squared norm of the basis function is
\begin{equation}
    \int_0^bJ_0(k_mr)^2\,r\,dr = \frac{b^2}{2}\Big[J_0(k_mb)^2+J_1(k_mb)^2\Big] = \frac{b^2}{2}J_0(k_mb)^2\text{,}
    \label{eq:bessel-norm}
\end{equation}
whence the weight $2/(bJ_0(k_mb))^2$ used above. \cite{AlventosaEtAl2023}'s eq.~(300) instead prints $2/(bJ_1(k_m))^2$, which is the normalizer for a \emph{Dirichlet} basis ($J_0(k_mb)=0$, $k_mb$ a zero of $J_0$) and moreover evaluates $J_1$ at $k_m$ rather than $k_mb$. Under the no-flux condition they themselves impose ($\partial_n\phi=0$ on the walls, hence $\partial_r\eta=0$ at $r=b$), $J_1(k_mb)$ vanishes identically, so their expression is inconsistent with their own basis. The discrepancy is mode-dependent and locally severe rather than a uniform rescaling: evaluated at $b=6$ against \eqref{eq:bessel-norm}, the ratio of their weight to the correct one runs $0.57,2.67,5.35,6.41,4.74,1.68,0.0086,1.28$ for $m=1,\dots,8$, the near-vanishing entry at $m=7$ arising because $k_7=3.7933$ falls close to the first zero of $J_1$, driving $J_1(k_m)\to0$. An earlier revision of this document reproduced their expression verbatim; every bath-side quantity reported here uses \eqref{eq:bessel-norm}.

\subsection{Droplet interface model}
\label{subsec:drop}

The droplet's small-amplitude oscillations about its spherical equilibrium are governed, for a general contact pressure $p(\theta,\tau)$, by the energy balance $\dot T=\dot W-\varepsilon$ of \cite{AlventosaEtAl2023} (their eq.~3.7), with $T$ the sum of kinetic and surface energy of the decomposition \eqref{eq:eta-xi}, $\dot W$ the rate of working of $p$ on the surface, and $\varepsilon$ the viscous dissipation of \cite{lamb1924hydrodynamics}. Substituting \eqref{eq:eta-xi} and using the orthogonality of the $P_l(\cos\theta)$ on $[-1,1]$ gives, mode by mode,
\begin{equation}
    \ddot\beta_l + 2\Oh(2l+1)(l-1)\, \dot\beta_l + l(l-1)(l+2)\, \beta_l = -(2l+1)l\, b_l(\tau)\text{,}\qquad l=2,\dots,L\text{,}
    \label{eq:drop-mode}
\end{equation}
where
\begin{equation}
    b_l(\tau) = \int_{-1}^1 p(x,\tau)\, P_l(x)\, dx
    \label{eq:b_l-def}
\end{equation}
is the un-normalized Legendre projection of the actual pressure, in the convention \cite{AlventosaEtAl2023} use for the projection $c_l$ of their prescribed shape; here $b_l$ plays the role their product $f(\tau)c_l(\tau)$ played.

The two parent models print forcings for \eqref{eq:drop-mode} that differ by $(2l+1)^2$, so a remark on which is adopted, and why, is warranted. \cite{AgueroEtAl2026} write the forcing as $-l\mathcal B_l$ with $\mathcal B_l$ the coefficient in $p=\sum_l\mathcal B_lP_l(\cos\theta)$, but state its inverse transform as $\mathcal B_l=(2l+1)^{-1}\int p\,P_l\,dx$; since $\int_{-1}^1P_l^2\,dx=2/(2l+1)$, the inverse of their own expansion is $\mathcal B_l=\tfrac12(2l+1)\int p\,P_l\,dx$, larger by $(2l+1)^2/2$. Their mode equation is correct under the latter reading, and an independent derivation settles which: for $\xi=R+\sum_lA_lP_l$ the kinetic and surface energies are $2\pi\rho R^3\dot A_l^2/(l(2l+1))$ and $2\pi\sigma(l-1)(l+2)A_l^2/(2l+1)$, and the rate of working of $p$ is $-2\pi R^2\dot A_lb_l$, so $\tfrac{d}{d\tau}[T+E_s]=\dot W$ gives $\rho R^3\ddot A_l=-\sigma l(l-1)(l+2)A_l-\tfrac12R^2l(2l+1)b_l$. This reproduces \cite{AlventosaEtAl2023}'s dimensional equation coefficient for coefficient, and matches \cite{AgueroEtAl2026}'s only with $\mathcal B_l=\tfrac12(2l+1)b_l$, confirming that the discrepancy is a misprinted projection formula in the latter rather than a difference of physics --- in particular it is not a lab-frame versus drop-frame difference, since both papers refer the deformation modes to the droplet's own centre of mass, and a change of frame is a rigid translation that projects entirely onto $l=1$. The residual factor of two between \eqref{eq:drop-mode} and the derivation above is the pressure scale: \eqref{eq:drop-mode} and \eqref{eq:com} are jointly consistent with $2\sigma/R$, matching \cite{AlventosaEtAl2023}'s force scale $2\pi\sigma R$ and their normalization $f=2\int_0^{r_c}p\,r\,dr$. Modes $l=0$ (excluded by incompressibility of the droplet) and $l=1$ do not appear in \eqref{eq:eta-xi}; $l=1$ is instead governed by the motion of the centre of mass,
\begin{equation}
    \dot z_{cm} = v\text{,}\qquad \dot v = \tfrac32 f(\tau) - \Bo\text{,}\qquad f(\tau) = 2\int_0^{r_c(\tau)} p(r,\tau)\, r\, dr\text{,}
    \label{eq:com}
\end{equation}
with $r_c(\tau)$ the (as yet undetermined) contact radius. Under the ansatz $p=fH_r(r/r_c)$ of \cite{AlventosaEtAl2023}, whose normalization $\int_0^b H_r(r/r_c)\,r\,dr=\tfrac12$ they impose, $2\int_0^{r_c}p\,r\,dr=2f\int_0^{r_c}H_r(r/r_c)\,r\,dr=2f\cdot\tfrac12=f$, so \eqref{eq:com} reduces to their own definition of the net contact force $f$ when the shape is prescribed, and generalizes it when it is not.

\subsection{Time discretization: affine response to pressure}
\label{subsec:affine}

\hspace*{1em}Equations \eqref{eq:bath-mode}, \eqref{eq:drop-mode} and the pair \eqref{eq:com} are each, mode by mode, of the form
\begin{equation}
    \ddot x + 2\gamma\dot x+\omega^2x = Fu(\tau)\text{,}
    \label{eq:generic-ode}
\end{equation}
with $x$ one of $a_m,\beta_l,z_{cm}$, $u$ the corresponding pressure moment $c_m,b_l,f$, and $(\gamma,\omega^2,F)$ read off from \eqref{eq:bath-mode}, \eqref{eq:drop-mode}, \eqref{eq:com} respectively:
\begin{equation}
    (\gamma,\omega^2,F)_{\text{bath}} = \Big(2\Oh k_m^2,\ (k_m^2+\Bo)k_m\tanh(k_mh_0),\ -2k_m\tanh(k_mh_0)\Big)\text{,}
    \label{eq:bath-coeffs}
\end{equation}
\begin{equation}
    (\gamma,\omega^2,F)_{\text{drop}} = \Big(\Oh(2l+1)(l-1),\ l(l-1)(l+2),\ -(2l+1)l\Big)\text{,}\qquad
    (\gamma,\omega^2,F)_{cm} = \left(0,\ 0,\ \tfrac32\right)\text{.}
    \label{eq:drop-com-coeffs}
\end{equation}
We discretize \eqref{eq:generic-ode} in time with the variable-step BDF2 scheme of \cite{AgueroEtAl2026}: with $s^k=\delta_t^k/\delta_t^{k-1}$, $\mathfrak a^k=(1+2s^k)/(1+s^k)$, $b^k=-(1+s^k)$, $c^k=(s^k)^2/(1+s^k)$, and $\dot f^{k+1}\approx(\mathfrak a^kf^{k+1}+b^kf^k+c^kf^{k-1})/\delta_t^k$ for any variable $f$, the first-order system $(x,y{=}\dot x)$ becomes
\begin{equation}
    \mathfrak a x^{k+1}+bx^k+cx^{k-1} = \delta\, y^{k+1}\text{,}\qquad
    \mathfrak a y^{k+1}+by^k+cy^{k-1} = \delta\left(-2\gamma y^{k+1}-\omega^2x^{k+1}+Fu^{k+1}\right)\text{,}
    \label{eq:bdf2}
\end{equation}
where $\delta:=\delta_t^k$ and the step superscript on $\mathfrak a,b,c$ is dropped for brevity. Solving the first equation of \eqref{eq:bdf2} for $y^{k+1}=(\mathfrak a x^{k+1}+bx^k+cx^{k-1})/\delta$ and substituting into the second gives
\begin{equation}
    \frac{\mathfrak a}{\delta}\Big(\mathfrak a x^{k+1}+bx^k+cx^{k-1}\Big)+by^k+cy^{k-1}
    = -2\gamma\Big(\mathfrak a x^{k+1}+bx^k+cx^{k-1}\Big)-\delta\omega^2x^{k+1}+\delta Fu^{k+1}\text{,}
\end{equation}
and collecting every term proportional to $x^{k+1}$ on the left,
\begin{equation}
    x^{k+1}\left[\frac{\mathfrak a^2}{\delta}+2\gamma\mathfrak a+\delta\omega^2\right]
    = \delta Fu^{k+1} - \frac{\mathfrak a}{\delta}\big(bx^k+cx^{k-1}\big)-2\gamma\big(bx^k+cx^{k-1}\big)-by^k-cy^{k-1}\text{,}
\end{equation}
so that, multiplying through by $\delta$,
\begin{equation}
    x^{k+1} = (\text{known history}) + \kappa\,u^{k+1}\text{,}\qquad
    \kappa = \frac{\delta^2 F}{\mathfrak a^2+2\gamma\mathfrak a\delta+\delta^2\omega^2}\text{.}
    \label{eq:kappa-general}
\end{equation}
Substituting \eqref{eq:bath-coeffs}--\eqref{eq:drop-com-coeffs} into \eqref{eq:kappa-general} gives, explicitly,
\begin{equation}
    a_m^{k+1} = \alpha_m^{k+1} + \kappa_m\, c_m^{k+1}\text{,}\qquad
    \kappa_m = \frac{-2\delta^2 k_m\tanh(k_mh_0)}{\mathfrak a\big(\mathfrak a+4\delta\Oh k_m^2\big)+\delta^2(k_m^2+\Bo)k_m\tanh(k_mh_0)}\text{,}
    \label{eq:kappa-m}
\end{equation}
\begin{equation}
    \beta_l^{k+1} = \gamma_l^{k+1} + \lambda_l\, b_l^{k+1}\text{,}\qquad
    \lambda_l = \frac{-\delta^2(2l+1)l}{\mathfrak a\big(\mathfrak a+2\delta\Oh(2l+1)(l-1)\big)+\delta^2l(l-1)(l+2)}\text{,}
    \label{eq:lambda-l}
\end{equation}
\begin{equation}
    z_{cm}^{k+1} = \mu^{k+1}+\kappa_{cm}\,f^{k+1}\text{,}\qquad \kappa_{cm}=\frac{3\delta^2}{2\mathfrak a^2}\text{,}
    \label{eq:kappa-cm}
\end{equation}
with $\alpha_m^{k+1},\gamma_l^{k+1},\mu^{k+1}$ depending only on $\delta,\mathfrak a,b,c$ and the two preceding time levels. Each of $a_m^{k+1}$, $\beta_l^{k+1}$, $z_{cm}^{k+1}$ is therefore an affine function of the pressure moment forcing it at the same time level; the slopes $\kappa_m,\lambda_l,\kappa_{cm}$ are known once $\delta,\mathfrak a$ are fixed by the time step, and it is this affine dependence that closes the contact problem below into a purely algebraic one at each time step.

\subsection{Modeling contact}
\label{subsec:contact}

\hspace*{1em}We now determine the pressure and the contact angle $\theta_c(\tau)$ so that the bath and droplet surfaces coincide throughout $[0,\theta_c(\tau)]$.

\subsubsection{Pressure representation}
\label{subsubsec:pressure-representation}
We represent the pressure as a polynomial of degree $N$ in $x=\cos\theta$, compactly supported on the pressed region by construction. Two representations are available and both are viable. A nonzero polynomial (or any convergent real-analytic series) cannot vanish identically on a sub-interval without vanishing identically everywhere --- a degree-$N$ polynomial has at most $N$ roots, so forcing it to zero across a continuum forces every coefficient to zero --- so a representation over the full domain $x\in[-1,1]$ with ``zero outside contact'' imposed as a separate weak condition can only approximate that property, never satisfy it exactly. An earlier revision of this document treated that observation as disqualifying, and concluded that piecewise support was a \emph{structural requirement} of a genuinely spectral pressure unknown. That conclusion was wrong, and is withdrawn: weak enforcement of $p\equiv0$ outside contact has exactly the standing every other Galerkin condition in this document has, including \eqref{eq:galerkin} itself, which likewise cannot enforce its own pointwise condition exactly. Rejecting the global representation on grounds that condemn the rest of the formulation equally was a double standard, not an argument.

We nonetheless adopt the piecewise representation, for a reason that survives scrutiny: with the contact angle determined by the outer scalar iteration of \S\ref{subsubsec:contact-angle} rather than carried as an unknown alongside the pressure coefficients, the support endpoint $x_c$ is a fixed \emph{parameter} throughout each inner solve. Exact vanishing outside contact is then free --- it costs no extra equations, introduces no approximation, and leaves nothing to be checked afterward --- so there is no reason to trade it for the weaker global alternative.

Within that constraint, the basis functions themselves are not required to be monomials, and monomials are in fact a poor choice: writing $\psi(x,\tau):=(x-x_c(\tau))/(1-x_c(\tau))\in[0,1]$ for the pressure's own rescaled coordinate, the monomials $\psi^n$ have a self-overlap (Gram) matrix $\int_0^1\psi^n\psi^m\,d\psi=1/(n+m+1)$ --- the classical Hilbert matrix, whose condition number grows like $1.6\times10^4$ already at $N=3$, independent of $x_c$. We instead use the shifted Legendre polynomials $\tilde P_n(\psi):=P_n(2\psi-1)$, exactly orthogonal under the same plain-$L^2$ inner product used throughout this document (the same one Gauss--Legendre quadrature and the droplet's own $P_l(x)$ basis already use, unlike Chebyshev polynomials, whose optimality is with respect to a different, weighted inner product that does not transfer to the mixed projections computed below), giving
\begin{equation}
    p(x,\tau) = \sum_{n=0}^{N} \hat c_n(\tau)\, \tilde P_n\big(\psi(x,\tau)\big)\text{,}\quad x\ge x_c(\tau)\text{;}\qquad
    p(x,\tau) \equiv 0\text{,}\quad x<x_c(\tau)\text{,}\qquad x_c(\tau):=\cos\theta_c(\tau)\text{.}
    \label{eq:pressure-poly}
\end{equation}
$\tilde P_n(\psi(x,\tau))$ is, for every $n$, exactly a degree-$n$ polynomial in $x$ (an affine change of variable composed with a degree-$n$ polynomial), so \eqref{eq:pressure-poly} has exactly the same total polynomial degree $N$ as the monomial basis it replaces --- confirmed symbolically, not merely asserted, in \texttt{derivations/verify\_legendre\_pressure\_basis.jl}. Their self-Gram matrix is, by orthogonality, diagonal: $\mathrm{cond}=2N+1$ at every $N$ and every $x_c$, against $1.6\times10^4$ for the monomial basis at $N=3$.

\subsubsection{Full contact}
Contact over the entire pressed region requires
\begin{equation}
    \eta\big(r(\theta,\tau),\tau\big) = z_{cm}(\tau) - z_d(\theta,\tau)\text{,}\qquad 0\le\theta\le\theta_c(\tau)\text{,}
    \label{eq:full-contact}
\end{equation}
with $\eta$, $z_{cm}$, $z_d$ affine in $c_m$, $f$, $b_l$ respectively by \eqref{eq:kappa-m}--\eqref{eq:kappa-cm}. The $N+1$ coefficients $\hat c_n(\tau)$ cannot satisfy \eqref{eq:full-contact} at every $\theta\in[0,\theta_c(\tau)]$; we impose it in the weighted-residual sense obtained by projecting the mismatch onto the $N+1$ polynomials used in \eqref{eq:pressure-poly},
\begin{equation}
    \int_{x_c}^1 \Big[\eta\big(r(\theta,\tau),\tau\big) - z_{cm}(\tau) + z_d(\theta,\tau)\Big]\,\tilde P_n\big(\psi(x,\tau)\big)\, w(x,\tau)\, dx = 0\text{,}\qquad n=0,\dots,N\text{,}
    \label{eq:galerkin}
\end{equation}
$N+1$ equations. By \eqref{eq:lambda-l}, $\beta_l(\tau)=\gamma_l(\tau)+\lambda_lb_l(\tau)$, and $b_l(\tau)=\int_{x_c}^1p(x,\tau)P_l(x)\,dx$ depends on $\hat c_n(\tau)$ alone --- not on $\eta$, $r$, or $z_d$ --- so $\beta_l(\tau)$ is an explicit function of the same unknowns \eqref{eq:galerkin} is solving for. Substituting it makes \eqref{eq:galerkin} genuinely nonlinear in $\hat c_n(\tau)$: $\beta_l(\tau)$ enters $\xi$, hence $r(\theta,\tau)$, hence $J_0\big(k_mr(\theta,\tau)\big)$ inside $\eta$ (through $c_m(\tau)$, \S\ref{subsubsec:quadrature} below), and $J_0$ composed with an affine function of $\hat c_n(\tau)$ is not itself affine. \eqref{eq:galerkin} is solved for $\hat c_n(\tau)$ at a fixed trial contact angle by the inner nonlinear solve of \S\ref{subsec:newton}, with $\theta_c(\tau)$ itself fixed by the outer scalar selection of \S\ref{subsubsec:contact-angle}; we first evaluate the four projections entering \eqref{eq:galerkin} and \eqref{eq:c_m-def} through $z_d$, $z_{cm}$ and $\eta$, each as an explicit function of $\hat c_n(\tau)$ and $\theta_c(\tau)$, in turn.

\subsubsection{The polynomial projections}
\label{subsubsec:polynomial-projections}
The centre-of-mass and droplet-surface contributions to \eqref{eq:com}, \eqref{eq:b_l-def} and \eqref{eq:galerkin} involve only $\xi$ and $p$, hence only polynomials in $x$, and are therefore \emph{exactly} (not merely approximately) evaluable by Gauss--Legendre quadrature \eqref{eq:gauss-quad}, given enough nodes: an $n_q$-point rule is exact for any polynomial of degree at most $2n_q-1$, and every integrand here has an explicitly known, finite degree. This replaces a closed-form antiderivative recursion (viable only for the monomial basis \eqref{eq:pressure-poly} previously used here, via the Legendre antiderivatives of \cite{AgueroEtAl2026}'s App.~A) with the same quadrature machinery already required below for the bath's $J_0$ terms, which have no antiderivative in closed form regardless of basis --- one uniform mechanism for every projection in this document, not two.

$b_l(\tau)=\int_{x_c}^1p(x,\tau)P_l(x)\,dx$ of \eqref{eq:b_l-def} has integrand $\tilde P_n(\psi(x,\tau))P_l(x)$, degree $N+L$ in $x$ (degree $N$ from \eqref{eq:pressure-poly}, degree $L$ from $P_l$), hence exact for $n_q\ge\lceil(N+L+1)/2\rceil$ quadrature points. The droplet contribution to \eqref{eq:galerkin},
\begin{equation}
    \int_{x_c}^1z_d(\theta,\tau)\,\tilde P_n\big(\psi(x,\tau)\big)\,dx
    = \int_{x_c}^1x\Big[1+\sum_{l=2}^L\beta_l(\tau)P_l(x)\Big]\tilde P_n\big(\psi(x,\tau)\big)\,dx\text{,}
    \label{eq:drop-galerkin-term}
\end{equation}
has integrand degree $N+L+1$ (the extra factor of $x$ in $z_d=\xi\cos\theta$), exact for $n_q\ge\lceil(N+L+2)/2\rceil$. Both degree counts are confirmed symbolically, not hand-counted, in \texttt{derivations/verify\_legendre\_pressure\_basis.jl}, together with the degree-$n$ claim for $\tilde P_n(\psi(x,\tau))$ itself.

The centre-of-mass integral in \eqref{eq:com} is likewise a pure polynomial integral: since $r(x,\tau)^2=\xi(x,\tau)^2(1-x^2)$ is itself a polynomial in $x$ (the $\sqrt{1-x^2}$ in $\sin\theta$ cancels once squared), $2\int_0^{r_c}p\,r\,dr=-\int_{x_c}^1p(x,\tau)\,d\big[r(x,\tau)^2\big]$ has integrand degree $N+2L+1$, exact for $n_q\ge\lceil(N+2L+2)/2\rceil$ --- unchanged from the monomial basis, since $p(x,\tau)$ still has total degree $N$ regardless of which degree-$N$ polynomial family represents it. The minus sign is not optional: $\theta=0$ (the pole, $r=0$) corresponds to $x=1$, and $\theta=\theta_c$ (the contact edge, $r=r_c$) corresponds to $x=x_c<1$, so $r$ \emph{increases} as $x$ \emph{decreases} over the contact patch; substituting $r\to x$ in $\int_0^{r_c}(\cdot)\,d[r^2]$ and reorienting the limits from $(1,x_c)$ to the natural $(x_c,1)$ introduces exactly this sign. Omitting it makes the pressure force accelerate the droplet further into the bath rather than decelerating it — confirmed by evaluating both sides directly for a constant test pressure $p\equiv1$, where the unsigned form returns $-r_c^2$ against the correct $2\int_0^{r_c}r\,dr=+r_c^2$.

\subsubsection{The quadrature projections}
\label{subsubsec:quadrature}
The bath-side projection $c_m$ of \eqref{eq:c_m-def} and the bath contribution to \eqref{eq:galerkin} both involve $J_0\big(k_mr(x,\tau)\big)$, and $J_0$ composed with the non-polynomial map $r(x,\tau)=\xi(\arccos x,\tau)\sqrt{1-x^2}$ has no elementary antiderivative. We evaluate
\begin{equation}
    c_m(\tau) = \frac{2}{(bJ_0(k_mb))^2}\int_{x_c}^1 p(x,\tau)\, J_0\!\big(k_mr(x,\tau)\big)\, r(x,\tau)\,\Big|\frac{dr}{dx}\Big|\, dx\text{,}\qquad
    W_n^{(m)}(\tau) := \int_{x_c}^1 J_0\!\big(k_mr(x,\tau)\big)\,\tilde P_n\big(\psi(x,\tau)\big)\,dx
    \label{eq:c_m-Wn}
\end{equation}
by $n_q$-point Gauss--Legendre quadrature on $[x_c,1]$: with $\{s_i,w_i\}_{i=1}^{n_q}$ the standard nodes and weights on $[-1,1]$ and the affine map $x=x_c+(1+s)(1-x_c)/2$,
\begin{equation}
    \int_{x_c}^1 g(x)\,dx \approx \frac{1-x_c}{2}\sum_{i=1}^{n_q} w_i\, g\!\left(x_c+\frac{(1+s_i)(1-x_c)}{2}\right)\text{,}
    \label{eq:gauss-quad}
\end{equation}
applied to $g$ equal to the integrand of $c_m$ and, separately, of $W_n^{(m)}$. No root is found in evaluating \eqref{eq:gauss-quad}: $r(x,\tau)$ is a forward evaluation of \eqref{eq:forward-map} at each quadrature abscissa, not an inversion, and no unknown of the problem is attached to an abscissa; \eqref{eq:c_m-Wn} is one definite integral per $(m,n)$, feeding linearly into \eqref{eq:bath-mode} and \eqref{eq:galerkin} respectively, in exactly the sense that any spectral method evaluates a forcing term that is not itself expressible in the trial basis by quadrature. With \eqref{eq:c_m-Wn}, the bath contribution to \eqref{eq:galerkin} is $\sum_m a_m(\tau)\,W_n^{(m)}(\tau)$.

\subsubsection{The compliance operator, and what the contact conditions determine}
\label{subsubsec:compliance}
\hspace*{1em}Before closing the system it is worth asking what object the closure is inverting, because its structure decides both what is determinable and how well conditioned the determination is. Collect the pressure-dependent part of the contact mismatch into a single linear map. By \eqref{eq:kappa-m}--\eqref{eq:kappa-cm} each of $\eta,z_{cm},z_d$ responds affinely to the pressure moments, and by \eqref{eq:c_m-def}, \eqref{eq:b_l-def} and \eqref{eq:com} each moment is a linear functional of $p$; composing, at frozen geometry,
\begin{equation}
    C(x,\tau) = C_{\text{free}}(x,\tau) + \big(\mathcal A p\big)(x,\tau)\text{,}\qquad
    \big(\mathcal A p\big)(x) = \int_{x_c}^1 \mathcal K(x,x')\,p(x')\,dx'\text{,}
    \label{eq:compliance-def}
\end{equation}
with $C_{\text{free}}$ the mismatch the BDF2 history alone would produce and $\mathcal K$ assembled from three contributions,
\begin{equation}
    \mathcal K = \underbrace{\sum_{m}\kappa_m\tfrac{2}{(bJ_0(k_mb))^2}J_0\big(k_mr(x)\big)J_0\big(k_mr(x')\big)\,w(x')}_{\text{bath}}
    \;\underbrace{-\,2\kappa_{cm}\,w(x')}_{\text{centre of mass}}
    \;+\;\underbrace{\sum_l\lambda_l\,x\,P_l(x)P_l(x')}_{\text{droplet}}\text{,}
    \label{eq:compliance-kernel}
\end{equation}
where $w(x):=-\tfrac12\,d[r^2]/dx$, which equals $r|dr/dx|$ and hence supplies the cylindrical area element $r\,dr=w\,dx$ \emph{provided} $dr/dx\le0$ across the patch --- a genuine restriction on the admissible contact extent, not an identity, recorded as such in \S\ref{subsubsec:admissibility}. $\mathcal A$ is the \emph{compliance operator} of the contact problem: it maps a trial pressure to the gap displacement that pressure induces. Its properties, measured in \texttt{derivations/audit\_compliance\_operator.jl} and \texttt{derivations/audit\_weak\_determinacy.jl}, govern what follows, and several of them are less favourable than an earlier revision of this section claimed.

\emph{$\mathcal A$ is self-adjoint under a consistent choice of pairing, and that choice is adopted.} Self-adjointness of $\mathcal A$ with respect to $\langle u,v\rangle_\mu=\int uv\,\mu\,dx$ requires $\mu(x)\mathcal K(x,x')=\mu(x')\mathcal K(x',x)$. Read off \eqref{eq:compliance-kernel}: the bath and centre-of-mass terms carry their entire $x'$-dependence in the factor $w(x')$, so they are self-adjoint precisely for $\mu\propto w$ --- confirmed at relative Frobenius asymmetry $\sim10^{-16}$ at every $\theta_c$ tested. The droplet term as written, $\sum_l\lambda_l\,xP_l(x)P_l(x')$, is not: it would require $\mu\propto1/x$ instead, and no weight satisfies both.

The obstruction is not, however, a property of the physics; it is a bookkeeping inconsistency, and it is the superposition of exactly two mismatches. First, $c_m$ and $f$ pair the pressure against the cylindrical element $w\,dx=r\,dr$ while $b_l$ \eqref{eq:b_l-def} pairs it against the spherical element $dx=\sin\theta\,d\theta$. Second, $b_l$ pairs against the droplet's \emph{radial} displacement, following \cite{AlventosaEtAl2023}'s energy balance, while the same mode returns to the contact mismatch through its \emph{vertical} component $x\xi$. Repairing either alone leaves the asymmetry non-zero; repairing \emph{both},
\begin{equation}
    b_l(\tau) \;=\; \int_{x_c}^1 p(x,\tau)\,x\,P_l(x)\,w(x,\tau)\,dx\text{,}
    \label{eq:b_l-selfadjoint}
\end{equation}
makes the droplet block --- and hence the whole of $\mathcal A$ --- self-adjoint in $\langle\cdot,\cdot\rangle_w$ to $10^{-16}$ at every deformation tested ($\beta_2=0,0.02,0.05,0.10$), against $3.6$--$5.5\times10^{-3}$ for \eqref{eq:b_l-def}. Both mismatches must be repaired together: taking the vertical factor alone leaves $1.8$--$3.6\times10^{-3}$ and the measure alone leaves $1.7$--$1.8\times10^{-3}$. An earlier revision of this section tested them only in isolation, concluded that neither was ``the'' mechanism, and inferred that self-adjointness was unattainable; that inference was wrong and is withdrawn (\S\ref{subsec:corrections}).

Equation \eqref{eq:b_l-selfadjoint} is adopted in place of \eqref{eq:b_l-def}. It is a change of \emph{linearisation convention}, not of physics, and it is $O(\theta_c^2)$: the pressure acts on the droplet surface along its true normal, which is neither exactly radial nor exactly vertical, and the two conventions differ only at second order in the contact angle. Three points make the choice the better one. It is internally consistent, treating the droplet's displacement in the same direction and against the same measure as the contact condition and the bath and centre-of-mass channels already use. It buys exact energy reciprocity of the coupled system, hence an energy functional and the monotone Galerkin convergence estimates that accompany one. And it degenerates to the parents' convention exactly in the limit the parents work in: $x\to1$ and $w\to1$ as $\theta_c\to0$, so \eqref{eq:b_l-selfadjoint} reduces to \eqref{eq:b_l-def} at pole contact, which is precisely where \cite{AlventosaEtAl2023} enforce their single-point match. Only \cite{AgueroEtAl2026} carry the inconsistency over a finite patch; \cite{AlventosaEtAl2023} do not incur it at all, so no appeal to inheritance is needed to justify departing from it.

One consequence is used below. With \eqref{eq:b_l-selfadjoint} in force, and with the contact-region Galerkin conditions tested against the \emph{same} basis that represents the pressure and in the \emph{same} measure $w\,dx$, the inner Galerkin matrix is symmetric --- measured at $\|G-G^{\!\top}\|/\|G+G^{\!\top}\|=2.1,1.9,1.3\times10^{-16}$ at $\theta_c=0.10,0.20,0.40$. Both conditions are necessary: pairing a different test basis against the trial basis destroys the symmetry even when the operator itself is self-adjoint, which is how the effect was first missed.

\emph{The symmetric part is positive semidefinite, not definite, and marginally so at larger patches.} Measured eigenvalues of the symmetrised $-\mathcal A$ in the $w$-measure, $n_q=40$: at $\theta_c=0.05,0.15,0.30,0.60$ the ratio $\lambda_{\min}/\lambda_{\max}$ is $-2.3\times10^{-16},-1.2\times10^{-16},-1.9\times10^{-12},-4.2\times10^{-7}$, with $35,31,26,18$ of $40$ eigenvalues below $10^{-12}\lambda_{\max}$. So the operator is semidefinite with a nullspace occupying roughly half the discretisation or more, and at $\theta_c=0.60$ the most negative eigenvalue is too large relative to $\lambda_{\max}$ to attribute to the $h=10^{-7}$ finite difference used for $d[r^2]/dx$ --- semidefiniteness is not cleanly established there. What semidefiniteness does buy is uniqueness of the gap response $\mathcal Ap$; it does \emph{not} buy uniqueness of $p$, and an earlier revision claimed the latter. Nor is a Frobenius-norm comparison of symmetric and antisymmetric parts evidence of definiteness in either direction, as that revision also supposed. This is not in tension with the compactness discussed below: a large numerical nullspace is precisely the finite-dimensional face of an operator whose singular values accumulate at zero.

\emph{The $O(\delta^2)$ pressure sensitivity is a near-uniform scalar factor only in a restricted regime.} Every term of \eqref{eq:compliance-kernel} carries exactly one factor of $\kappa_m$, $\kappa_{cm}$ or $\lambda_l$, each $O(\delta^2)$ by \eqref{eq:kappa-general}, which suggests $\mathcal A=\delta^2\widehat{\mathcal A}$ with $\widehat{\mathcal A}$ independent of $\delta$. But the denominators of \eqref{eq:kappa-m}--\eqref{eq:lambda-l} carry $\delta$-dependent, \emph{mode-dependent} terms, so the factorisation is only asymptotic, and the regime in which it holds is narrower than an earlier revision of this section asserted. The relative deviation of $\kappa_m/\delta^2$ from its $\delta\to0$ limit, maximised over $m$, is $0.011$ at $M=40,\delta=10^{-3}$, but $0.273$ at $M=160$ and $0.960$ at $M=700$; for $\lambda_l$ the corresponding figures are $0.052$, $0.691$ and $0.994$. Uniformity therefore requires $\delta k_M^2\ll1$ and $\delta^2L^3\ll1$, and is \emph{violated} at the mode counts used in the resolution study below. Correspondingly $\|\mathcal A\|/\delta^2$ is measured at $5.97\times10^4,1.51\times10^5,1.60\times10^5,1.60\times10^5$ across $\delta=10^{-2},\dots,10^{-5}$ --- settling, but varying by a factor of $2.5$ at the largest step, which the earlier revision suppressed by quoting only the limit. The important point is that the refutation of the deleted acceleration-level closure (\S\ref{subsec:corrections}) never needed uniformity: it needs only $\mathrm{cond}(c\mathcal A)=\mathrm{cond}(\mathcal A)$ together with the \emph{measured} $\delta$-insensitivity of the inner system, which is flat and of order ten (\S\ref{subsubsec:contact-angle}) and is a direct measurement rather than an inference.

\emph{$\mathcal A$ is compact, and its resolvable dimension must be measured rather than predicted.} $\mathcal A$ is an integral operator with a smooth kernel, hence compact, hence has singular values accumulating at zero. The evidence for this is the decay of the \emph{resolved} part of the spectrum, not a computed condition number: at $\sim10^{18}$ the latter exceeds $1/\epsilon_{\text{mach}}\approx4.5\times10^{15}$, so it reports that $\sigma_{\min}$ is unresolvable and carries no further information --- an earlier revision quoted its variation across $\delta$ as evidence of flatness, which it cannot be. What the resolved spectrum does show is that $\sigma_i/\sigma_1$ converges pointwise as $M,L$ grow, so $\mathcal A$ is a genuine limit object rather than an artifact of truncation, with the leading ratios still drifting at $M=L=160$ (for $i=2$: $0.352,0.318,0.272,0.254$ at $M=L=20,40,80,160$) while the tail fills in from below; and that the converged decay is algebraic rather than geometric, evidenced at large $i$ ($\sigma_{14}/\sigma_2=0.14\approx2/14$ at $M=L=160$, against $0.78^{12}=0.05$ for the geometric rate the first few ratios alone would also fit).

An earlier revision of this section reported a law $\mathrm{rank}\,\mathcal A\approx0.75\,k_Mr_c+2$ and attributed the resolvable pressure order to the bath truncation. That law is withdrawn. Its supporting study set $M=L$ throughout, in which $k_Mr_c$ and the droplet-side $L\theta_c$ are collinear, so it had no power to distinguish the two blocks; and it held $n_q$ fixed, so its largest cases were bounded by $\mathrm{rank}\le n_q$ rather than by the operator. With $M$ and $L$ varied independently and $n_q$ refined to independence, \emph{both} truncations are found to contribute, and which one binds depends on the state. Holding $M$ fixed and starving the droplet: at $L=160$, $\theta_c=0.4$ the rank is $25$ whether $M$ is $10$, $40$ or $160$ --- a sixteenfold change in $k_Mr_c$ with no effect --- while at $M=160$, $\theta_c=0.4$ it falls to $15$ for both $L=40$ and $L=10$. But holding $L$ fixed and enriching the bath: at $L=10$, $\theta_c=0.4$ the rank runs $5,6,15,46$ for $M=10,40,160,640$, and $46$ exceeds anything the droplet block can supply there ($L-1=9$ modes plus the rank-one centre-of-mass term), so the bath is contributing the remaining $36$ dimensions. $M$ is irrelevant only where the droplet term already saturates. Nor is $L\theta_c$ a substitute for $k_Mr_c$: $(M,L,\theta_c)=(160,40,0.4)$ and $(80,80,0.2)$ share $L\theta_c=16$ and give ranks $15$ and $9$. No single-parameter collapse survives the decoupling, and none is asserted here; the withdrawn law failed by omitting the droplet term, not because the bath term is absent. The measured rank is in addition bounded above by $n_q$, so it is reported only where refining $n_q$ (to $80,160,240$) leaves it unchanged, and it is $\delta$-independent (measured $14,15,15,15$ across $\delta=10^{-2},10^{-3},10^{-5},10^{-7}$). The practical prescription does not require one: the resolvable dimension is a property of the assembled discrete operator, obtained from its singular values at negligible cost relative to a time step, and it varies through an impact as the patch grows. It is therefore \emph{measured} at each step and used to set the spectral filter of \S\ref{subsec:newton}, rather than predicted from a formula.

\emph{The pointwise pressure does not converge; whether its moments do is not established here.} That the pointwise pressure is not a converged output is robust across every state and mode count tested: with the filter cutoff at $10^{-10}$ and $M=L=800$, $\theta_c=0.16$, $\min_xp$ runs $+98,\,-2.1\times10^3,\,-1.0\times10^6,\,-5.0\times10^5$ at $N=10,20,40,80$, and similarly at $M=L=400$ and at $\theta_c=0.40$. This is the expected behaviour of an ill-posed first-kind inversion and is not by itself a defect: only the moments $c_m,b_l,f$ enter \eqref{eq:bath-mode}, \eqref{eq:drop-mode} and \eqref{eq:com}, so only they need converge. An earlier revision claimed that they do. The evidence does not support that claim, and it is withdrawn. Over the same $N$ sequence at $M=L=800$, $\theta_c=0.16$, the relative change in $b_l$ between successive $N$ is $1.5\times10^{-2},3.6\times10^{-2},6.4\times10^{-5}$ --- non-monotone --- while at $\theta_c=0.40$ it is $1.9\times10^{-2},3.3\times10^{-2},4.4\times10^{-2}$, increasing rather than decreasing. Holding $N$ at $40$ and sweeping the filter cutoff over $10^{-6}$ to $10^{-14}$ moves $b_l$ by $\approx1.0\times10^{-2}$ per decade and $f$ from $11.28$ to $11.92$: the cutoff sensitivity is the same size as the $N$-drift, so the two cannot be separated by this data. And $f$ itself differs by $7\%$ between $M=400$ and $M=800$ at fixed $\theta_c=0.16$ ($12.55$ against $11.64$), so convergence in $M$ is not established either. The honest status is therefore that the closure's conditioning and structure are established (\S\ref{subsubsec:contact-angle}, \S\ref{subsec:newton}) while its convergence is not; establishing it requires a refinement study in $N$, $M$, $L$ and cutoff jointly, at admissible states, which is not carried out here.

\subsubsection{The contact angle: an outer selector, not a joint unknown}
\label{subsubsec:contact-angle}

\hspace*{1em}The single remaining unknown is $\theta_c(\tau)$. Write
\begin{equation}
    C(\theta,\tau) := \eta\big(r(\theta,\tau),\tau\big)-z_{cm}(\tau)+z_d(\theta,\tau)
    \label{eq:pointwise-residual-sec3}
\end{equation}
for the pointwise contact mismatch, which \eqref{eq:full-contact} requires to vanish for $0\le\theta\le\theta_c(\tau)$ and, by \eqref{eq:check-nonintersect} below, to remain negative beyond it.

\emph{Tangency, and the sense in which it is degenerate.} The classical closing condition for a free boundary of this kind is tangency --- smooth pasting between the pressed and unpressed regions, $\partial_\theta C(\theta_c,\tau)=0$. It carries an exact degeneracy at the poles. By \eqref{eq:eta-xi}, $\xi(\theta,\tau)$ depends on $\theta$ only through $\cos\theta$, hence is even in $\theta$ for \emph{every} $\theta$, not merely near the axis; $r(\theta,\tau)=\xi(\theta,\tau)\sin\theta$ is therefore odd; $\eta(r(\theta,\tau),\tau)=\sum_ma_m(\tau)J_0(k_mr)$ is even, since $J_0$ is an even function of its (here, odd-in-$\theta$) argument, so the composition is even; and $z_d(\theta,\tau)=\xi(\theta,\tau)\cos\theta$ is even$\times$even, hence even. Every formula composing $C(\theta,\tau)$ ($\cos\theta$, $\sin\theta$, $J_0$, the finite Legendre sum in $\xi$) is entire and defined for every real $\theta$, so this parity statement is a genuine global identity about the analytic extension of $C$ to all of $\mathbb R$, not an assumption smuggled in by extending the physical domain. $C$ is therefore even in $\theta$ at every $\theta$ and every physical state, its $\theta$-derivative is odd and $2\pi$-periodic, and an odd $2\pi$-periodic function vanishes at both fixed points of $\theta\mapsto-\theta$:
\begin{equation}
    \partial_\theta C(0,\tau)\equiv0\text{,}\qquad \partial_\theta C(\pi,\tau)\equiv0
    \label{eq:tangency-degenerate}
\end{equation}
identically, for \emph{any} $a_m,\beta_l$ whatsoever. Near $\theta=0$, $\partial_\theta C(\theta)=C''(0)\theta+O(\theta^3)$ is affine through the origin for any state, so a Newton iteration on $\partial_\theta C=0$ seeded anywhere near onset finds the trivial root essentially independent of the physics --- confirmed directly: joint solves seeded at $\theta_c^{(0)}\in\{0.02,\dots,0.5\}$, at various levels of accumulated contact history, all converged to $\theta_c\sim10^{-8}$.

An earlier revision of this document read \eqref{eq:tangency-degenerate} as showing tangency to be \emph{structurally} unusable, and discarded it. That inference was too strong, and is withdrawn (\S\ref{subsec:corrections}). Equation \eqref{eq:tangency-degenerate} is a statement about the poles only, and it is the correct \emph{physical} degeneracy there: at $\theta_c=0$ the contact patch has shrunk to a point and tangency carries no information about its extent, because there is no extent. The identical degeneracy holds in \cite{AgueroEtAl2026}'s own formulation --- axisymmetry forces $\partial_r\eta(0,t)=\partial_rz_d(0,t)=0$ for every state, so their tangency condition \citep{AgueroEtAl2026} is likewise vacuous at $r_c=0$ --- and they use it successfully throughout an impact regardless. The difference is not in the condition but in how it is used: they never solve it by a root-finder seeded at an arbitrary contact extent. They evaluate its residual at a handful of candidate contact extents bracketing the previous step's, reject those violating non-intersection outright, and take the argument minimising the residual among the survivors. Tangency is serviceable as a \emph{selector} over candidate contact extents; it is degenerate only as an \emph{equation solved from a free seed}. What follows adopts the former role.

\emph{The nesting, and why it is not circular.} Let $\theta_c$ be an outer scalar \emph{parameter} rather than an unknown carried alongside the pressure. For each trial value, $x_c=\cos\theta_c$ is fixed, the pressure representation \eqref{eq:pressure-poly} and every projection of \S\ref{subsubsec:polynomial-projections}--\S\ref{subsubsec:quadrature} are well defined, and the $N+1$ Galerkin conditions \eqref{eq:galerkin} form a square system for the $N+1$ coefficients $\hat c_n$ alone. Solving it yields a pressure, hence --- through \eqref{eq:summary-states} --- a complete state, hence a scalar residual in $\theta_c$ with the pressure already eliminated. The outer problem is one-dimensional.

An earlier revision asserted that computing $\theta_c$ this way is circular, on the grounds that $p(x;X)$'s support depends on $\theta_c$, so $b_l,c_m,a_m,\beta_l$ cannot be evaluated without $\theta_c$ already known. That argument is correct, and it rules out the \emph{opposite} nesting --- computing $\theta_c$ from an already-determined $\hat c$, downstream of quantities that presuppose it. It says nothing against the nesting above, in which $\theta_c$ is supplied \emph{first}, as a parameter, and everything downstream is evaluated from it in the same directed acyclic order \S\ref{subsec:newton} already establishes. Conflating the two directions was the error; \S\ref{subsec:corrections} records it.

\emph{The selectors, and their agreement.} Two scalar conditions are available to fix $\theta_c$ once the pressure is eliminated, and they are not independent. The first is tangency in its selector role,
\begin{equation}
    T(\theta_c,\tau) := \partial_\theta C(\theta,\tau)\big|_{\theta=\theta_c}\text{,}\qquad\text{sought }=0\text{,}
    \label{eq:tangency-selector}
\end{equation}
evaluated with the state induced by the pressure that \eqref{eq:galerkin} returns at that same $\theta_c$. The second is the non-intersection condition \eqref{eq:check-nonintersect}, imposed not as an afterthought but as a feasibility filter on candidate $\theta_c$ exactly as \cite{AgueroEtAl2026} impose it: any trial value for which $C(\theta,\tau)\ge0$ for some $\theta>\theta_c$ is inadmissible and discarded before $T$ is consulted at all.

How well these two locate the same free boundary is a numerical question, and the answer is weaker than an earlier revision of this section claimed. At the representative penetrating state of \texttt{derivations/audit\_nested\_closure.jl} ($M=L=80$, $\delta=10^{-3}$, centre of mass at $\mu=0.985$, so the undeformed surfaces cross at $\arccos\mu=0.1734$), scanning $\theta_c$ on a grid of step $0.005$ gives
\begin{center}
\begin{tabular}{lccc}
 & $T=0$ & \eqref{eq:check-nonintersect} becomes feasible & $\min_xp$ changes sign \\
$N=1$ & $0.1925$ & $0.1725$ & $0.1975$ \\
$N=2$ & $0.1925$ & $0.1675$ & $0.2025$ \\
$N=3$ & $0.1925$ & $0.1625$ & $0.1125$
\end{tabular}
\end{center}
The three do \emph{not} coincide --- they spread over $0.025$ to $0.080$ depending on $N$, and the positivity crossing is erratic at $N=3$ --- so the earlier claim that they agree on one interval, and the inference that the closure was thereby exhibiting a complementarity structure, are both withdrawn (\S\ref{subsec:corrections}). What the table does show is more useful for the closure's actual purpose: the tangency root is $0.1925$ at \emph{every} $N$, identically to four figures, so the selector is insensitive to the pressure truncation; and the non-intersection boundary sits within $1\%$ of the parameter-free geometric estimate $\arccos\mu$, with the tangency root about $11\%$ above it. Pressure positivity is not imposed anywhere in this closure --- consistent with both parent models, neither of which constrains the sign of $p$ pointwise (\cite{AlventosaEtAl2023} monitor only $\mathrm{sign}\,f$; \cite{AgueroEtAl2026}'s constraint list contains no positivity condition) --- and it holds on roughly the same interval without being asked to.

\emph{The selector is an argmin over candidates, not a bisection, and not an integrated residual.} $\theta_c^{k+1}$ is fixed exactly as \cite{AgueroEtAl2026} fix their contact-point count: evaluate the residual at a finite set of candidate values bracketing the previous step's converged $\theta_c^k$, discard those failing the feasibility filters \eqref{eq:check-nonintersect} and \eqref{eq:check-monotone-r}, and take
\begin{equation}
    \theta_c^{k+1} \;=\; \operatorname*{arg\,min}_{\theta\in\Theta^{k+1}}\big|T(\theta,\tau^{k+1})\big|\text{,}
    \label{eq:theta-c-argmin}
\end{equation}
$\Theta^{k+1}$ the admissible candidates. Two properties of this construction are not incidental. It needs no sign change, so the measured non-monotonicity of $T$ --- which does occur, and which can place two roots inside one bracket --- does not defeat it, whereas a bisection on $T$ would require a bracket guaranteed to contain exactly one root and no such guarantee is available. And $T$ is a \emph{pointwise} quantity evaluated at the contact edge, of fixed size as $\theta_c$ varies.

That last point is a genuine constraint on the design, established by trying the alternative. A selector minimising an \emph{integrated} weak residual over the contact patch --- for instance the least-squares residual of a system combining the kinematic conditions on $[x_c,1]$ with weakly-imposed zero pressure on $[-1,x_c]$, using a pressure basis global on $[-1,1]$ --- is biased toward vanishing contact, because shrinking the patch shrinks the domain the residual is measured over. Measured (\texttt{derivations/probe\_global\_basis\_argmin.jl}): that residual is monotone increasing away from $\theta_c=0$, so its argmin is the smallest admissible $\theta_c$, reproducing the onset degeneracy rather than curing it. The same probe records two further obstructions to the global-basis route: the number of test functions imposing zero pressure off the patch must be strictly less than $N+1$, since $N+1$ of them form a square full-rank system admitting only $p\equiv0$ (the identity-theorem obstruction in weak form); and with fewer, the combined system is numerically rank deficient, $\mathrm{cond}\sim10^{14}$ to $10^{22}$, so its least-squares residual sits at roundoff for every $\theta_c$ and carries no signal at all. \cite{AgueroEtAl2026}'s argmin is not subject to any of this precisely because their residual is a pointwise edge quantity; \eqref{eq:theta-c-argmin} follows them.

\emph{Conditioning.} Removing $\theta_c$ from the linear system is what repairs the conditioning, and the mechanism is now unambiguous. By \S\ref{subsubsec:compliance}, every column of the pressure block carries a factor of $\kappa_m$, $\kappa_{cm}$ or $\lambda_l$, all $O(\delta^2)$; a $\theta_c$ column does \emph{not}, since $\theta_c$ is a raw coordinate that never flows through the BDF2 relations \eqref{eq:kappa-m}--\eqref{eq:kappa-cm}. A joint system therefore mixes columns differing by $\delta^2$ in scale, and its condition number grows as $\delta$ shrinks --- exactly the measured behaviour of the joint formulation. The inner system of the nesting contains no such column, and its condition number is measured (\texttt{derivations/audit\_nested\_closure.jl}, $N=3$, $M=L=80$, $\theta_c=0.3$) at $8.19,\,3.11,\,2.97,\,2.96,\,2.96$ for $\delta=10^{-2},\dots,10^{-6}$: flat in $\delta$ and of order unity, against $\sim10^{18}$ for the joint system. The gain is roughly eighteen orders of magnitude and is obtained by removing one row, not by adding machinery.

\emph{Bracketing, in place of a multiplicity argument.} Because $\eta(r(\theta,\tau),\tau)$ is a sum of oscillatory $J_0$ terms once $k_mr$ is not small, $T(\cdot,\tau)$ may change sign more than once on $(0,\pi)$, and $C(\cdot,\tau)$ likewise. This is the same multi-rootedness a joint Newton solve would face, but it is far more tractable in one dimension: the outer problem admits bracketing. Following \cite{AgueroEtAl2026}, the search is confined to a neighbourhood of the previous step's converged $\theta_c^k$ and the accepted value is required to satisfy a continuity bound $|\theta_c^{k+1}-\theta_c^k|$ below a threshold, with the time step reduced and the step retried when it is not --- a physical requirement (the contact line cannot jump) rather than a numerical expedient, and the direct analogue of their $|r_c^{k+1}-r_c^k|\le\delta_r$. Within such a bracket \eqref{eq:theta-c-argmin} selects among finitely many candidates, which cannot return a value outside the bracket.

\emph{Onset --- not specified here.} One step is not covered by the construction above, and the gap is stated rather than papered over. At first contact $\theta_c^k=0$, and \eqref{eq:tangency-degenerate} makes $T\equiv0$ identically there for every state, so $T$ is uninformative in a neighbourhood of $\theta_c=0$ and the first bracket cannot be centred on the previous step's value. The scanned data show the degeneracy reaching into the working range rather than staying at the endpoint: $T$ is non-monotone and decreasing toward the bottom of the admissible window. \cite{AgueroEtAl2026} guard this explicitly --- their discrete tangency condition is stated ``if $q>0$'', and their candidate loop runs over integers so that $q=0$ is a distinguishable, excluded state --- and this document has no analogue of either guard. A pragmatic initialisation is available and is what the implementation uses: at the first contact step, seed $\theta_c$ from the geometric crossing $C(\theta_c,\tau)=0$, which is non-degenerate at onset (\S\ref{subsubsec:contact-angle}'s existence argument) and which the table above places within $1\%$ of the reference, then hand over to \eqref{eq:theta-c-argmin} from the second step on. That is an initialisation choice, not a derived rule, and it is recorded as such.

Existence of at least one admissible $\theta_c$ for a state genuinely in contact follows as before from the intermediate value theorem applied to $C$: $C(0,\tau)>0$ is the geometric onset-of-contact condition this closure presupposes (the pole's absolute height is $z_{cm}-\xi(0,\tau)$, so the gap there is $-C(0,\tau)$, negative exactly when the pole has penetrated), while $C(\theta,\tau)$ must become negative before $\theta=\pi$, since a positive $C$ all the way to $\pi$ would mean the surfaces overlap everywhere --- not a state this document's kinematic-match premise admits. Continuity of $C$ in $\theta$ then gives a crossing. This is a plausibility check on the closure, not a guarantee that the bracketed search finds the physically continuous branch at every step; that remains a numerical matter, monitored as described.

\subsubsection{Admissibility}
\label{subsubsec:admissibility}
Two conditions are structural in the closure above, and one is a diagnostic rather than a constraint. Non-intersection of the two surfaces beyond the contact edge,
\begin{equation}
    \eta\big(r(\theta,\tau),\tau\big) < z_{cm}(\tau)-z_d(\theta,\tau)\text{,}\qquad \theta\in(\theta_c(\tau),\pi]\text{,}
    \label{eq:check-nonintersect}
\end{equation}
is a feasibility filter on candidate $\theta_c$, applied before the selector \eqref{eq:tangency-selector} is consulted, exactly as in \cite{AgueroEtAl2026}. The kinematic match itself holds in the weighted-residual sense \eqref{eq:galerkin} on the contact patch, and $p\equiv0$ off it holds exactly by construction of \eqref{eq:pressure-poly}.

A second condition is structural and was left implicit in an earlier revision. Every substitution of $r\,dr=w\,dx$ in \eqref{eq:c_m-Wn}, \eqref{eq:compliance-kernel} and \eqref{eq:summary-moments} requires $x\mapsto r(x,\tau)$ to be injective on the patch, equivalently
\begin{equation}
    \frac{\partial r}{\partial x}(x,\tau) < 0\text{,}\qquad x\in[x_c(\tau),1]\text{,}
    \label{eq:check-monotone-r}
\end{equation}
without which $w=-\tfrac12d[r^2]/dx$ is not $r|dr/dx|$, can change sign, and the change of variables is invalid. For an undeformed droplet $r=\sqrt{1-x^2}$ and \eqref{eq:check-monotone-r} holds throughout; for a deformed one it fails once the patch reaches the widest point of the droplet, which is exactly \cite{AgueroEtAl2026}'s $r_M(t)$ --- ``the smallest distance from the $z$ axis to the point on the surface of the droplet for which the plane that is tangent to the interface of the droplet is vertical'', which they introduce precisely as an upper bound on the admissible contact radius. Condition \eqref{eq:check-monotone-r} is the same bound, $r_c(\tau)<r_M(\tau)$, and must be tested on each candidate $\theta_c$ alongside \eqref{eq:check-nonintersect}.

Positivity of the pressure throughout the pressed region,
\begin{equation}
    p(x,\tau) = \sum_{n=0}^N\hat c_n(\tau)\,\tilde P_n\big(\psi(x,\tau)\big) \ge 0\text{,}\qquad x\in[x_c(\tau),1]\text{,}
    \label{eq:check-positivity}
\end{equation}
is neither imposed nor, at large $N$, a meaningful test. It is not imposed because neither parent model imposes it and because it is observed to hold on the admissible interval without being asked to (\S\ref{subsubsec:contact-angle}). It ceases to be a meaningful test at large $N$ for the reason established in \S\ref{subsubsec:compliance}: the compliance operator is compact, so the pointwise pressure is not a converged output of the model at any $M$, and $\min_xp$ is measured to diverge with $N$. A pointwise sign test applied to a quantity the model determines only weakly reports the truncation, not the physics. At small $N$, where the pressure is close to fully resolved, \eqref{eq:check-positivity} remains informative and is worth monitoring; at larger $N$ the corresponding diagnostic is sensitivity of $c_m,b_l,f$ to the spectral-filter cutoff, since it is the cutoff and not $N$ that sets the resolution there (\S\ref{subsec:newton}). Neither diagnostic is a convergence proof: \S\ref{subsubsec:compliance} records that convergence of the moments is not established.

Nor does \eqref{eq:pressure-poly} constrain the value of the pressure at the edge of contact: $p(x_c,\tau)=\hat c_0(\tau)-\hat c_1(\tau)+\cdots$ is generally nonzero, admitting a jump in $p$ across $\theta_c(\tau)$. This is not obviously a defect. \cite{AlventosaEtAl2023}'s review of direct numerical simulation and experiment on this problem (not solid--solid Hertzian contact) finds the true pressure distribution during impact closer to flat than singular at the edge, which is why they fit a high-order rather than singular shape function --- so a prescribed vanishing at $\theta_c$ is not evidently the right target, and the edge behaviour is a property of the solution rather than a constraint to impose in advance.

\subsection{Solution procedure}
\label{subsec:newton}

\hspace*{1em}The contact problem at time level $\tau^{k+1}$ is solved as a one-dimensional outer iteration over $\theta_c^{k+1}$ wrapping an inner solve for the pressure coefficients. The inner unknowns are
\begin{equation}
    \hat c := \big(\hat c_0^{k+1},\dots,\hat c_N^{k+1}\big)\in\mathbb R^{N+1}\text{,}
    \label{eq:newton-unknowns}
\end{equation}
with $\theta_c^{k+1}$ --- and hence $x_c^{k+1}=\cos\theta_c^{k+1}$ --- held fixed at the current trial value throughout. Every other quantity is an explicit function of $\hat c$, computed in the following order with no step depending on a quantity not yet available: $p(x;\hat c)$ from \eqref{eq:pressure-poly}, using the fixed $x_c^{k+1}$; $b_l(\hat c)$ and the droplet contribution to \eqref{eq:galerkin} from \eqref{eq:b_l-def} and \eqref{eq:drop-galerkin-term}; $\beta_l(\hat c)=\gamma_l^{k+1}+\lambda_lb_l(\hat c)$ from \eqref{eq:lambda-l}; $\xi(x;\hat c)=1+\sum_l\beta_l(\hat c)P_l(x)$, $r(x;\hat c)=\xi\sqrt{1-x^2}$, $z_d(x;\hat c)=\xi x$ from \eqref{eq:eta-xi}--\eqref{eq:forward-map}; $c_m(\hat c)$ and $W_n^{(m)}(\hat c)$ from \eqref{eq:c_m-Wn}; $a_m(\hat c)=\alpha_m^{k+1}+\kappa_mc_m(\hat c)$ from \eqref{eq:kappa-m}; and $f(\hat c)$, $z_{cm}(\hat c)=\mu^{k+1}+\kappa_{cm}f(\hat c)$ from \eqref{eq:com}, \eqref{eq:kappa-cm}. The inner residual is
\begin{equation}
    R_n(\hat c) := \int_{x_c}^1\Big[\eta\big(r(\theta;\hat c),\tau^{k+1}\big)-z_{cm}(\hat c)+z_d(\theta;\hat c)\Big]\,\tilde P_n\big(\psi(x)\big)\,dx\text{,}\quad n=0,\dots,N\text{,}
    \label{eq:newton-residual-galerkin}
\end{equation}
$N+1$ equations for $N+1$ unknowns, solved by Newton's method
\begin{equation}
    \hat c^{(j+1)} = \hat c^{(j)} - J\big(\hat c^{(j)}\big)^{-1}R\big(\hat c^{(j)}\big)\text{,}\qquad J:=\frac{\partial R}{\partial\hat c}\text{.}
    \label{eq:newton-step}
\end{equation}
The system is genuinely nonlinear, though only mildly: by \S\ref{subsubsec:compliance} it would be exactly linear, $R=\langle\tilde P_n,C_{\text{free}}+\mathcal A\hat c\rangle$, were the geometry frozen; the nonlinearity enters entirely through the $\beta_l(\hat c)$-dependence of $r(x;\hat c)$: non-polynomially inside $J_0(k_mr)$, and \emph{quadratically} through $r^2=\xi^2(1-x^2)$, which appears in $f$ and in $w$ and therefore makes $z_{cm}(\hat c)$ a quadratic rather than affine function of $\hat c$. Both channels are $O(\delta^2)$-suppressed by \eqref{eq:lambda-l}, which is why Newton converges rapidly from the previous time level's converged $\hat c$ --- an expectation, not yet a measurement. What has been measured is the conditioning of the frozen-geometry Galerkin matrix $\langle\tilde P_n,\mathcal A\tilde P_{n'}\rangle$ that $J$ reduces to when the geometry is held fixed: of order unity and flat in $\delta$ (\S\ref{subsubsec:contact-angle}), provided $N+1$ does not exceed the measured resolvable dimension of \S\ref{subsubsec:compliance}. Neither $\mathrm{cond}(J)$ at a converged nonlinear state nor the Newton iteration count has been measured here. Because every arrow in the chain above (polynomial evaluation, Gauss--Legendre quadrature \eqref{eq:gauss-quad}, and $J_0,J_1$ with their derivative recurrences $J_0'=-J_1$, $J_1'=J_0-J_1/x$) is an elementary or special-function composition, and none is a linear solve, an implicit root or a search, $R$ is differentiable by direct automatic differentiation of the same computation that evaluates it; forming $J$ costs $N+1$ forward-mode passes and requires no adjoint rule.

Equation \eqref{eq:newton-step}'s linear solve is always performed as a spectral-filtered solve --- a truncated singular value decomposition of $J$ at a fixed relative cutoff, discarding directions below it --- and not only when $N$ is large. This is the standard treatment of a compact operator's ill-posed inversion rather than a tuning knob in disguise: by \S\ref{subsubsec:compliance} the discarded directions are those in which the pressure does not measurably affect the gap, and the retained count is exactly the measured resolvable dimension, which adapts on its own as the patch grows through the impact. Adopting it unconditionally is what makes the policy consistent: the resolution is set by the cutoff, so $N$ may be over-provisioned freely, and no rule for choosing $N$ correctly is needed. The price is that the cutoff is now the parameter the solution depends on, and \S\ref{subsubsec:compliance} records that this dependence is of the same size as the residual $N$-dependence, so it must be reported rather than assumed negligible. Choosing $N$ strictly inside the resolvable dimension, so that no filtering is active, remains available and is the regime in which $\hat c(\theta_c)$ is unambiguously single-valued.

The outer iteration is then a one-dimensional search for $\theta_c^{k+1}$: bracket around $\theta_c^k$ subject to the continuity bound, discard trial values violating \eqref{eq:check-nonintersect}, and bisect on the sign change of \eqref{eq:tangency-selector} among the survivors, reducing $\delta$ and retrying the step if no admissible bracket exists. Each outer evaluation costs one inner solve. The scheme is fully implicit in the pressure and requires no extrapolated state anywhere.

\subsection{Model summary}
\label{subsec:summary}

The complete system solved at every time level $\tau^{k+1}$: the $N+1$ equations \eqref{eq:summary-galerkin} for $\hat c^{k+1}$ at fixed $\theta_c^{k+1}$, wrapped in the scalar selection \eqref{eq:summary-theta-c} for $\theta_c^{k+1}$ itself, together with the states, geometry and pressure moments that \S\ref{subsec:newton} shows to be explicit functions of $\hat c^{k+1}$ alone.

\emph{States, advanced from \eqref{eq:kappa-m}--\eqref{eq:kappa-cm} given the pressure moments at $\tau^{k+1}$:}
\begin{equation}
    a_m^{k+1}=\alpha_m^{k+1}+\kappa_mc_m^{k+1}\ (m=0,\dots,M)\text{,}\quad
    \beta_l^{k+1}=\gamma_l^{k+1}+\lambda_lb_l^{k+1}\ (l=2,\dots,L)\text{,}\quad
    z_{cm}^{k+1}=\mu^{k+1}+\kappa_{cm}f^{k+1}\text{.}
    \label{eq:summary-states}
\end{equation}

\emph{Geometry, given the state:}
\begin{equation}
    r(\theta,\tau^{k+1})=\xi(\theta,\tau^{k+1})\sin\theta\text{,}\qquad z_d(\theta,\tau^{k+1})=\xi(\theta,\tau^{k+1})\cos\theta\text{,}\qquad
    \xi=1+\textstyle\sum_l\beta_l^{k+1}P_l(\cos\theta)\text{.}
    \label{eq:summary-geometry}
\end{equation}

\emph{Pressure, the inner unknowns, at the current trial contact angle:}
\begin{equation}
    p(x,\tau^{k+1})=\sum_{n=0}^N\hat c_n^{k+1}\,\tilde P_n\big(\psi(x,\tau^{k+1})\big)\ (x\ge x_c^{k+1})\text{,}\qquad p\equiv0\ (x<x_c^{k+1})\text{,}\qquad x_c^{k+1}=\cos\theta_c^{k+1}\text{.}
    \label{eq:summary-pressure}
\end{equation}

\emph{Pressure moments feeding \eqref{eq:summary-states}, computed from \eqref{eq:summary-pressure} via \eqref{eq:b_l-def}, \eqref{eq:drop-galerkin-term} and \eqref{eq:c_m-Wn}, all by quadrature \eqref{eq:gauss-quad}:}
\begin{equation}
    c_m^{k+1}=\tfrac{2}{(bJ_0(k_mb))^2}\!\int_{x_c^{k+1}}^1\! p\,J_0(k_mr)\,r\Big|\tfrac{dr}{dx}\Big|dx\text{,}\quad
    b_l^{k+1}=\!\int_{x_c^{k+1}}^1\! p\,x\,P_l(x)\,w\,dx\text{,}\quad
    f^{k+1}=-\!\int_{x_c^{k+1}}^1\! p\,d\big[r^2\big]\text{.}
    \label{eq:summary-moments}
\end{equation}

\emph{Inner closure, $N+1$ equations for the $N+1$ unknowns $\hat c_0^{k+1},\dots,\hat c_N^{k+1}$ at fixed $\theta_c^{k+1}$:}
\begin{equation}
    \int_{x_c^{k+1}}^1\Big[\eta\big(r(\theta,\tau^{k+1}),\tau^{k+1}\big)-z_{cm}^{k+1}+z_d(\theta,\tau^{k+1})\Big]\,\tilde P_n\big(\psi(x,\tau^{k+1})\big)\,w\,dx=0\text{,}\quad n=0,\dots,N\text{,}
    \label{eq:summary-galerkin}
\end{equation}

\emph{Outer closure, selecting $\theta_c^{k+1}$ among candidates admissible under \eqref{eq:check-nonintersect}, \eqref{eq:check-monotone-r} and the continuity bound:}
\begin{equation}
    \theta_c^{k+1}=\operatorname*{arg\,min}_{\theta\in\Theta^{k+1}}\big|\partial_\theta C(\theta,\tau^{k+1})\big|\text{,}\qquad
    C(\theta,\tau) = \eta\big(r(\theta,\tau),\tau\big)-z_{cm}(\tau)+z_d(\theta,\tau)\text{.}
    \label{eq:summary-theta-c}
\end{equation}

The pressure and the contact extent are both solved for --- neither is prescribed --- but they are solved at different levels rather than jointly, which is what keeps the inner system's conditioning independent of $\delta$ and of order unity. $\eta$, $\xi$ and $p$ remain finite spectral sums throughout; \eqref{eq:forward-map} is used only forward; $p\equiv0$ off the contact patch holds exactly; and the resolvable pressure order is measured from the assembled compliance operator at each step (\S\ref{subsubsec:compliance}) rather than chosen or predicted.

\section{Validation}
\label{sec:validation}

\hspace*{1em}The closure of \S\ref{sec:spectral-km} has been implemented and run against
published data for the water case of \cite{AlventosaEtAl2023} ($\We=1.0958$, $\Bo=0.017$,
$\Oh=0.006$), using the experimental and direct-numerical-simulation records that
accompany \cite{AgueroEtAl2026} rather than \cite{AlventosaEtAl2023}'s own reduced
curves, since the former supersede the latter. At $M=L=60$, $N=3$, $n_q=40$,
$\delta_{\text{init}}=10^{-3}$:

\begin{center}
\begin{tabular}{lccccc}
series & $n$ & mean $|$err$|$ & max $|$err$|$ & RMS & RMS/bar \\
droplet top & 60 & $0.1345$ & $0.2503$ & $0.1526$ & $1.52$ \\
droplet bottom & 39 & $0.1055$ & $0.2580$ & $0.1221$ & $1.22$ \\
combined & 99 & $0.1231$ & $0.2580$ & $0.1414$ & $1.41$ \\
\end{tabular}
\end{center}

against a measurement uncertainty whose median half-bar is $0.1004$ (quartiles $0.0810$
and $0.1179$, from 293 digitised bar caps). The model tracks the measured trajectory at
about $1.4$ times the experimental error bar --- the same order as the measurement
scatter, but outside it, so the residual is model error rather than noise. The trajectory
is physically bounded and rebounds: the centre of mass enters contact at $v=-1.047$ and
leaves at $v=+0.314$.

What can \emph{not} be validated this way must be said plainly, because an earlier
revision of this section reported it as though it could. The dataset contains no measured
contact radius and no measured droplet width; both exist only as DNS and 1PKM series,
verified by reading the digitiser projects that generated them, each of which holds
exactly two datasets. Any statement about contact radius is therefore a comparison against
another computation rather than against measurement, and is additionally sensitive to
definition: the DNS contact radius is thresholded on the trapped gas film, whereas
$\theta_c$ here bounds the kinematically matched region, and the two need not coincide.

With that caveat the contact-radius behaviour is still worth recording. Against DNS the
model's $r_c$ peaks at $\tau\approx0.89$ against $1.46$, and refining $M$ and $L$ moves
that by only $0.02$, so it is not a resolution effect. But the early peak is partly a
property of this class of reduced model rather than of this model alone: 1PKM peaks at
$\tau\approx1.22$, also early against the same DNS. Normalised by contact time, the peak
falls at $0.20$ of the way through contact here, $0.26$ for 1PKM and $0.33$ for DNS --- so
this model is more skewed than 1PKM, in the same direction, with no measurement available
to adjudicate. Whether the remainder is a defect of the closure or a mismatch of
definitions is open.

Two discrepancies are larger than the rest and are recorded rather than explained. The
contact radius is $7\%$ low at its maximum, and --- more tellingly --- reaches that maximum
at $\tau=0.87$ against the DNS value $\tau=1.46$, so the modelled contact patch both grows
and recedes too quickly even though the total contact time is accurate to $3\%$. Whether
this is a resolution effect or a defect of the closure is not settled here; it is the
natural first target for the next iteration.

A resolution study sharpens \S\ref{subsubsec:compliance}'s weak-determinacy finding, and
does so on a real trajectory rather than the synthetic state used there. Sweeping the
pressure truncation at $M=L=60$, the contact time is $4.26239$ at every one of
$N=2,3,6,10,16$ --- identical to six significant figures --- while the solved pressure
profile at $N=3$ differs from that at $N=16$ by $46$ to $96\%$ in relative $L^2$ over the
patch, and the net force by between $10^{-5}$ and $10^{-2}$. Both statements are correct
and they are not in tension: the pressure acts on the dynamics only through $c_m$, $b_l$
and $f$, and the modes that distinguish $N=3$ from $N=16$ lie in the numerical nullspace
of $\mathcal A$. The practical consequence is that trajectories may be computed at $N=3$,
but the pointwise pressure must not be quoted as a converged output at any $N$ --- which
is the same conclusion \S\ref{subsubsec:compliance} reaches from the operator's spectrum,
now confirmed dynamically.

One observable does drift with $N$: the maximum contact radius falls monotonically from
$0.8192$ at $N=2$ to $0.7999$ at $N=16$, moving \emph{away} from the DNS value $0.8821$.
That the single quantity most sensitive to the pressure profile is also the one carrying
the open timing discrepancy is unlikely to be a coincidence, and is the most concrete lead
available for it. Convergence is judged here on contact time rather than the coefficient
of restitution deliberately: the latter settles far earlier than anything else and would
have reported convergence at mode counts where the contact time is still wrong by a factor
of two ($M=L=30$ gives $2.15$ against the converged $4.26$).

One implementation defect worth recording, because its symptom was easy to mistake for a
cost scaling rather than a bug. Loss of contact was originally declared as soon as the net
force $f$ turned non-positive, following \cite{AlventosaEtAl2023}'s own $\mathrm{sign}\,f$
criterion. That criterion is too eager here: $f$ can dip transiently negative during the
early transient without contact ending --- measured at $\We=0.3$, where it runs
$+0.0002,\,+0.0063,\,+0.0367,\,-0.0038,\,+0.0321,\,+0.0269,\dots$ --- and acting on the
single excursion returns the stepper to free flight, where onset is immediately
re-detected and a fresh onset search is paid. Repeated, the cycle cost that Weber number
some twenty times the wall clock of its neighbours ($555$ s against $23$--$32$ s), which
looked like a scaling pathology and was not. Requiring $f\le0$ on several consecutive
steps removes it: $\We=0.3$ now runs in $24$ s, and the contact time becomes monotone in
$\We$ across the range tested ($4.768,\,4.388,\,4.262,\,4.186$ at
$\We=0.0646,\,0.3,\,1.0958,\,3.0$), as it should be. The complementary failure is also
real and is guarded separately: as the patch collapses $f\to0$ from above and never
changes sign, so the primary loss-of-contact test is the contact angle reaching its search
floor, not the force.

The implementation exercises exactly the structure derived above, and two of its details
were forced by measurement rather than chosen. First, reducing the time step does
\emph{not} help a step whose candidates all fail the non-intersection filter: the
pressure's ability to deform either interface within a step scales as $\delta^2$
(\S\ref{subsubsec:compliance}), so a smaller step has strictly less authority to remove an
overlap, and the usual shrink-and-retry response is counterproductive here. Second, the
feasible set of contact angles is a \emph{band}, not a half-line --- infeasible below,
where the surfaces interpenetrate beyond the edge, and infeasible again above, where
\eqref{eq:check-monotone-r} fails once the patch reaches the droplet's widest point --- so
the lower edge must be bracketed by a search that starts from the previous step's angle
and walks outward. A plain bisection assuming monotonicity walks past the band and reports
no admissible angle at steps where one exists.

\section{Corrections to earlier revisions}
\label{subsec:corrections}

\hspace*{1em}This document has been revised substantially, and several arguments in earlier revisions were wrong rather than merely superseded. They are recorded here because each was used to justify structure that has now been removed, and because two of them were load-bearing for an entire section.

\emph{The $O(\delta^2)$ conditioning diagnosis, and the acceleration-level closure built on it.} An earlier revision contained a section arguing that the position-level closure \eqref{eq:galerkin} \emph{could not} be well conditioned, because $\eta,z_{cm},z_d$ respond to the current step's pressure only at $O(\delta^2)$; it proposed replacing \eqref{eq:galerkin} by a Galerkin projection of $\ddot C=0$, obtained by differentiating the contact condition twice in $\tau$ and substituting the governing ODEs, so that the pressure would enter at $O(1)$. The premise is false, though not for the reason a previous revision of \emph{this} section gave. That revision asserted the $O(\delta^2)$ factor to be exactly \emph{uniform} across the pressure block; \S\ref{subsubsec:compliance} now measures the factorisation $\mathcal A=\delta^2\widehat{\mathcal A}$ to be asymptotic only, failing badly ($96\%$ deviation in $\kappa_m/\delta^2$) at the mode counts of interest. The refutation does not need uniformity. It needs only that a scalar factor cannot change a condition number, $\mathrm{cond}(c\mathcal A)=\mathrm{cond}(\mathcal A)$, together with the direct measurement that the inner system's condition number is flat in $\delta$ and of order unity once $\theta_c$ is removed from it. The measured $\delta$-dependence of the joint system's Jacobian came from mixing the $O(\delta^2)$ pressure columns with a $\theta_c$ column carrying no such factor; removing $\theta_c$ from the linear system, as \S\ref{subsubsec:contact-angle} now does, leaves a condition number flat in $\delta$ and of order unity. The acceleration-level section has accordingly been deleted in full, along with the open ``hidden constraint drift'' problem it introduced --- an artifact of differentiating a constraint that did not need differentiating, and not a defect of the underlying model. The residual ill-conditioning that section attributed partly to the pressure basis is now identified as the compactness of $\mathcal A$ (\S\ref{subsubsec:compliance}): $\delta$-independent, basis-independent in origin, and correctly handled by spectral filtering rather than by raising the differential order of the closure.

\emph{Tangency as ``structurally unusable.''} The degeneracy \eqref{eq:tangency-degenerate} is real, but the inference that tangency cannot serve as a closing condition was too strong. It is a statement about the poles, where it is the correct physical degeneracy, and it holds identically in \cite{AgueroEtAl2026}, who use tangency throughout an impact nonetheless. What \eqref{eq:tangency-degenerate} rules out is solving tangency by a root-finder seeded at an arbitrary contact extent, not tangency itself. \S\ref{subsubsec:contact-angle} uses it in the selector role its degeneracy permits.

\emph{The circularity argument.} The claim that determining $\theta_c$ from the pressure is circular is correct for the nesting in which $\theta_c$ is computed \emph{after} $\hat c$, and was mistakenly applied to the nesting in which $\theta_c$ is supplied first as a parameter. The two are not the same, and the latter is not circular.

\emph{Piecewise pressure support as a ``structural requirement.''} The identity-theorem argument of \S\ref{subsubsec:pressure-representation} was used to reject a global pressure representation with weakly-imposed vanishing outside contact. The observation is true but not disqualifying, since it condemns \eqref{eq:galerkin} on the same grounds; the piecewise representation is retained on the different and weaker ground that the outer nesting makes it free.

\emph{$M$, $L$ and $N$ as free parameters, and a withdrawn rank law.} An earlier revision treated the pressure truncation $N$ as an independent choice and recorded as an unresolved concern that $\mathrm{cond}(J)$ grew sharply with it. A subsequent revision claimed to resolve this with a law $\mathrm{rank}\,\mathcal A\approx0.75\,k_Mr_c+2$, attributing the resolvable pressure order to the bath truncation. That law is withdrawn (\S\ref{subsubsec:compliance}): its supporting study set $M=L$, making the claimed collapse variable collinear with a droplet-side one, and held $n_q$ fixed, capping its largest cases at the quadrature count; decoupling $M$ from $L$ shows both truncations to contribute --- a ninefold rank swing at fixed $L$ --- with no single-parameter collapse surviving. The resolvable dimension is measured from the assembled operator instead. Relatedly, that revision's claim that the pressure's moments converge in $N$ while its pointwise values do not is also withdrawn: only the second half is supported, and the residual drift in the moments is the same size as the spectral filter's own cutoff sensitivity.

\emph{The self-adjointness diagnosis, twice wrong in opposite directions.} One revision of \S\ref{subsubsec:compliance} stated that the droplet block's asymmetry is caused by ``a single factor of $x$'', that repairing $b_l\to\int p\,xP_l\,dx$ removes it, and that the operator is monotone because its symmetric part dominates in Frobenius norm. All three were wrong: the repair was measured in a different inner product from the one the claim was stated in; the two candidate mechanisms are indistinguishable at $\beta=0$, where $w\equiv x$ exactly, so no measurement taken at the undeformed state could have separated them; and norm dominance of the symmetric part is not evidence of definiteness in either direction. The following revision over-corrected, concluding from the failure of each repair \emph{in isolation} that no mechanism was identifiable and that self-adjointness was unattainable. That is also withdrawn: repairing both mismatches together, \eqref{eq:b_l-selfadjoint}, makes the whole operator self-adjoint to $10^{-16}$, and the formulation now adopts it. What survives unchanged from the first correction is the definiteness statement: the operator is measured positive semidefinite with a large nullspace, giving uniqueness of the gap response and not of the pressure.

\emph{The three-selector agreement.} A revision reported that the tangency root, the non-intersection feasibility boundary and the pressure-positivity crossing all fall between $\theta_c=0.16$ and $0.18$ at every $N$, and read this as evidence of a complementarity structure. Those numbers predated the Bessel-normalization correction and do not survive it: on a refined grid the three spread over $0.025$ to $0.080$ (\S\ref{subsubsec:contact-angle}). The claim is withdrawn. What replaces it is narrower and more directly useful --- the tangency root is $N$-independent to four figures --- and the structural inference drawn from the coincidence is retracted with it.

\emph{Bisection, and the integrated-residual selector.} A revision specified the outer selector as a bisection on the sign change of $T$; another proposed replacing it with an argmin of an integrated weak residual over a global pressure basis. Both are withdrawn in favour of \eqref{eq:theta-c-argmin}. Bisection requires a bracket containing exactly one root, which the measured non-monotonicity of $T$ does not provide, and is undefined at onset. The integrated-residual selector is biased toward vanishing contact, because shrinking the patch shrinks the domain the residual is measured over; measured, its argmin is the smallest admissible $\theta_c$, so it reproduces the onset degeneracy instead of curing it, and the global basis it requires is separately obstructed (\S\ref{subsubsec:contact-angle}). The lesson, stated because it was not obvious in advance, is that \cite{AgueroEtAl2026}'s argmin works because its residual is a \emph{pointwise edge} quantity of fixed size, not because it is an argmin.
